\chapter{سامانه‌های توصیه‌گر \texorpdfstring{\enparen{Recommender Systems}}{(Recommender Systems)}}
\label{ch:recsys}

\section{مقدمه و پدیده دنباله طولانی}
در دنیای فیزیکی، فروشگاه‌ها (نظیر کتاب‌فروشی‌ها یا سینماها) به دلیل محدودیت فضای قفسه‌ها، ناگزیرند تنها محبوب‌ترین و پرفروش‌ترین کالاها را عرضه کنند (پدیده فرهنگ عامه‌پسند یا پرفروش‌ها\footnote{\lr{Blockbuster Culture}}). اما در اقتصاد دیجیتال، وب‌سایت‌هایی مانند آمازون، نتفلیکس یا اسپاتیفای دارای هزینه‌های انبارداری و تکثیر نزدیک به صفر هستند؛ در نتیجه می‌توانند میلیون‌ها کالای خاص و کم‌یاب را نیز ارائه نمایند.

پدیده \textbf{دنباله طولانی}\footnote{\lr{Long Tail}} نشان می‌دهد که اگرچه تقاضا برای هر کالای خاصِ موجود در انتهای دنباله اندک است، اما \textbf{مجموع تقاضای تمامی اقلام دنباله طولانی} می‌تواند از تقاضای کالاهای پرفروش سنتی فراتر رود. با این حال، در میان میلیون‌ها گزینه، کاربر چگونه کالای مورد علاقه خود را پیدا کند؟ اینجاست که \textbf{سامانه‌های توصیه‌گر}\footnote{\lr{Recommender Systems}} به عنوان موتور محرک اقتصاد وب ظاهر می‌شوند.

\section{ماتریس سودمندی \texorpdfstring{\enparen{Utility Matrix}}{(Utility Matrix)}}
ستون فقرات ریاضی هر سامانه توصیه‌گر، \textbf{ماتریس سودمندی}\footnote{\lr{Utility Matrix}} است:

\begin{definition}[ماتریس سودمندی]
ماتریس $M$ با ابعاد $|U| \times |I|$ تعریف می‌شود که در آن سطرها کاربران $(U)$، ستون‌ها آیتم‌ها یا کالاها $(I)$، و درایه‌های $r_{ui}$ بیانگر درجه علاقه یا امتیاز کاربر $u$ به آیتم $i$ هستند (مثلاً امتیازی بین ۱ تا ۵ ستاره).
\end{definition}

ویژگی بارز ماتریس سودمندی در کلان‌داده، \textbf{تنک بودن فوق‌العاده}\footnote{\lr{Extreme Sparsity}} آن است: یک کاربر معمولی ممکن است از میان ده‌ها هزار فیلم موجود در نتفلیکس، تنها چند ده فیلم را تماشا کرده و امتیاز داده باشد؛ بنابراین بیش از $99\%$ درایه‌های ماتریس خالی (مجهول) هستند.

\textbf{هدف بنیادین سامانه توصیه‌گر}: پیش‌بینی دقیق مقادیر درایه‌های خالی ماتریس سودمندی و پیشنهاد آیتم‌هایی که بالاترین امتیاز تخمینی را دارند به هر کاربر.

\section{روش‌های مبتنی بر محتوا \texorpdfstring{\enparen{Content-Based Recommendation}}{(Content-Based Recommendation)}}
در این رویکرد، توصیه بر اساس ویژگی‌های ذاتی کالاها و تاریخچه علایق گذشته همان کاربر انجام می‌پذیرد.

\subsection{پروفایل آیتم‌ها و پروفایل کاربر}
\begin{enumerate}
    \item \textbf{پروفایل آیتم \enparen{Item Profile}}: هر کالا با یک بردار ویژگی توصیف می‌شود. برای مثال، برای یک فیلم، بردار شامل ژانرها (اکشن، کمدی)، نام کارگردان، بازیگران، و وزن کلمات کلیدی خلاصه فیلم (به کمک معیار TF-IDF) است.
    \item \textbf{پروفایل کاربر \enparen{User Profile}}: سلیقه کاربر $u$ با ترکیب بردارهای ویژگی آیتم‌هایی که به آن‌ها امتیاز بالا داده است، به دست می‌آید:
    \begin{equation}
    w_u = \frac{\sum_{i \in I_u} r_{ui} \cdot x_i}{\sum_{i \in I_u} r_{ui}}
    \end{equation}
    که در آن $x_i$ بردار ویژگی آیتم $i$ و $r_{ui}$ امتیاز کاربر به آن است.
    \item \textbf{تخمین امتیاز}: میزان تناسب آیتم جدید $j$ برای کاربر $u$ با استفاده از \textbf{تشابه کسینوسی} میان بردار کاربر و بردار آیتم سنجیده می‌شود:
    \begin{equation}
    \text{Score}(u, j) = \cos(w_u, x_j) = \frac{w_u \cdot x_j}{\|w_u\| \|x_j\|}
    \end{equation}
\end{enumerate}

\begin{examnote}
مزایا و معایب فیلترینگ محتوامحور:
\begin{itemize}
    \item \textbf{مزیت}: نیازی به داده‌های سایر کاربران ندارد؛ برای آیتم‌های تازه اضافه شده به سیستم مشکل راه‌اندازی سرد ایجاد نمی‌کند.
    \item \textbf{عیب بزرگ}: \textbf{بیش‌تخصصی شدن}\footnote{\lr{Overspecialization}} و ایجاد «حباب فیلتر»؛ سامانه هرگز کالایی خارج از دایره واژگان یا سبک گذشته کاربر را به او پیشنهاد نمی‌دهد (فقدان شانس کشف غیرمنتظره یا Serendipity).
\end{itemize}
\end{examnote}

\section{فیلترینگ مشارکتی \texorpdfstring{\enparen{Collaborative Filtering}}{(Collaborative Filtering)}}
فیلترینگ مشارکتی نیازی به استخراج محتوا و ویژگی‌های متنی آیتم‌ها ندارد. در عوض، بر پایه این فرض عمل می‌کند: «افرادی که در گذشته سلیقه و امتیازدهی یکسانی داشته‌اند، در آینده نیز نظرات مشابهی خواهند داشت».

\subsection{کسینوس مرکزی (همبستگی پیرسون)}
برای مقایسه دو کاربر یا دو آیتم، نمی‌توان از تشابه کسینوسی معمولی استفاده کرد؛ زیرا کاربران سخت‌گیر میانگین نمراتشان پایین (مثلاً ۲ از ۵) و کاربران سهل‌گیر میانگین نمراتشان بالاست (مثلاً ۴ از ۵).
راهکار علمی، استفاده از \textbf{کسینوس مرکزی}\footnote{\lr{Centered Cosine}} یا ضریب همبستگی پیرسون است که در آن میانگین امتیازات هر کاربر از نمرات او کسر می‌شود.

\begin{example}[محاسبه کسینوس مرکزی برای دو کاربر]
فرض کنید امتیاز دو کاربر $A$ و $B$ به چهار فیلم به شرح زیر باشد (علامت $-$ به معنای عدم مشاهده است):
\begin{center}
\begin{tabular}{ccccc}
\toprule
کاربر & فیلم ۱ & فیلم ۲ & فیلم ۳ & فیلم ۴ \\
\midrule
$A$ & $4$ & $5$ & $3$ & $-$ \\
$B$ & $2$ & $3$ & $-$ & $5$ \\
\bottomrule
\end{tabular}
\end{center}
میانگین امتیازات کاربر $A$: $\bar{r}_A = (4 + 5 + 3)/3 = 4$.
میانگین امتیازات کاربر $B$: $\bar{r}_B = (2 + 3 + 5)/3 = 3.33$.

با نرمال‌سازی و کسر میانگین (تنها بر روی فیلم‌های مشترک ۱ و ۲):
\begin{align}
\tilde{A} &= [4-4, 5-4] = [0, 1] \\
\tilde{B} &= [2-3.33, 3-3.33] = [-1.33, -0.33]
\end{align}
کسینوس مرکزی:
\begin{equation}
\text{sim}(A, B) = \frac{0 \times (-1.33) + 1 \times (-0.33)}{\sqrt{0^2 + 1^2} \sqrt{(-1.33)^2 + (-0.33)^2}} = \frac{-0.33}{1 \times 1.37} \approx -0.24
\end{equation}
(نمره منفی نشان‌دهنده اختلاف سلیقه در فیلم‌های مشترک پس از حذف اثر سخت‌گیری است).
\end{example}

\subsection{پیش‌بینی امتیاز کاربر به آیتم (فیلترینگ کاربر-کاربر)}
برای پیش‌بینی امتیاز کاربر $u$ به آیتم $i$:
\begin{enumerate}
    \item مجموعه $N$ شامل $k$ کاربر شبیه‌ترین به $u$ را که به آیتم $i$ امتیاز داده‌اند، انتخاب می‌کنیم.
    \item امتیاز تخمینی با میانگین‌گیری وزنی بر اساس تشابه کسینوس مرکزی به دست می‌آید:
    \begin{equation}
    \hat{r}_{ui} = \bar{r}_u + \frac{\sum_{v \in N} \text{sim}(u, v) \cdot (r_{vi} - \bar{r}_v)}{\sum_{v \in N} |\text{sim}(u, v)|}
    \end{equation}
\end{enumerate}

\subsection{فیلترینگ آیتم-آیتم در برابر کاربر-کاربر}
در عمل (به‌ویژه در آمازون)، فیلترینگ \textbf{آیتم-آیتم} ترجیح داده می‌شود، زیرا:
\begin{itemize}
    \item ویژگی‌ها و ماهیت کالاها بسیار پایدارتر از سلیقه متغیر انسان‌ها در طول زمان است.
    \item شباهت میان آیتم‌ها را می‌توان به صورت آفلاین محاسبه و پیش‌محاسبه کرد؛ در نتیجه زمان پاسخگویی برخط بسیار سریع‌تر خواهد بود.
\end{itemize}

\section{مدل‌های عامل پنهان و فاکتورسازی ماتریسی \texorpdfstring{\enparen{Matrix Factorization}}{(Matrix Factorization)}}
پیشرفته‌ترین روش در فیلترینگ مشارکتی (که برنده جایزه یک میلیون دلاری نتفلیکس شد)، تجزیه ماتریس سودمندی به مؤلفه‌های پنهان\footnote{\lr{Latent Factors}} است.

\begin{definition}[تجزیه UV]
فرض کنید ماتریس سودمندی تنک $M$ دارای ابعاد $m \times n$ باشد. هدف، تقریب این ماتریس از حاصل‌ضرب دو ماتریس با رتبه پایین $d$ (که $d \ll \min(m, n)$ است، مثلاً $d = 20$ یا $50$) است:
\begin{equation}
M \approx P \cdot Q^T
\end{equation}
که در آن:
\begin{itemize}
    \item ماتریس $P$ با ابعاد $m \times d$، بردار عوامل پنهان کاربران است $(p_u)$.
    \item ماتریس $Q$ با ابعاد $n \times d$، بردار عوامل پنهان آیتم‌ها است $(q_i)$.
\end{itemize}
امتیاز تخمینی برابر است با ضرب داخلی: $\hat{r}_{ui} = p_u^T \cdot q_i$.
\end{definition}

\subsection{تابع هدف بهینه‌سازی و منظم‌سازی \texorpdfstring{\enparen{Regularization}}{(Regularization)}}
برای جلوگیری از بیش‌برازش بر روی داده‌های محدود مشاهده‌شده، تابع هزینه به صورت مجموع مربعات خطا به همراه ترم منظم‌سازی $L_2$ فرمول‌بندی می‌شود:
\begin{equation}
\min_{P, Q} \sum_{(u, i) \in \text{Observed}} (r_{ui} - p_u^T q_i)^2 + \lambda \left( \sum_u \|p_u\|_2^2 + \sum_i \|q_i\|_2^2 \right)
\end{equation}
که در آن $\lambda$ ضریب منظم‌سازی است.

\subsection{بهینه‌سازی با گرادیان کاهشی تصادفی \texorpdfstring{\enparen{SGD}}{(SGD)}}
برای هر درایه مشاهده‌شده $r_{ui}$، خطای پیش‌بینی برابر است با:
\begin{equation}
e_{ui} = r_{ui} - p_u^T q_i
\end{equation}
قاعده به‌روزرسانی پارامترها با نرخ یادگیری $\gamma$:
\begin{align}
p_u &\leftarrow p_u + \gamma (e_{ui} \cdot q_i - \lambda p_u) \\
q_i &\leftarrow q_i + \gamma (e_{ui} \cdot p_u - \lambda q_i)
\end{align}
این به‌روزرسانی بسیار ساده، سریع و با یک یا چند گذر روی کلان‌داده قابل اجرا است.

\subsection{مدل پایه با اثرات انحراف \texorpdfstring{\enparen{Baseline Predictors}}{(Baseline Predictors)}}
بسیاری از تغییرات امتیازات ناشی از سلیقه نیست، بلکه ناشی از رفتارهای سیستماتیک است:
\begin{equation}
\hat{r}_{ui} = \mu + b_u + b_i + p_u^T q_i
\end{equation}
که در آن:
\begin{itemize}
    \item $\mu$: میانگین کل امتیازات در سرتاسر پایگاه داده.
    \item $b_u$: انحراف کاربر $u$ از میانگین (سخت‌گیر یا سهل‌گیر بودن).
    \item $b_i$: انحراف کیفیت آیتم $i$ از میانگین (فیلمی فوق‌العاده یا ضعیف).
\end{itemize}

\begin{examnote}[مقایسه روش‌های بهینه‌سازی SGD و ALS در آزمون]
\begin{itemize}
    \item \textbf{گرادیان کاهشی تصادفی \enparen{SGD}}: در هر گام بر روی یک جفت کاربر-آیتم عمل می‌کند. پیاده‌سازی آن بسیار ساده و همگرایی اولیه آن سریع است، اما موازی‌سازی آن به دلیل هم‌پوشانی متغیرها در خوشه‌های توزیع‌شده چالش‌برانگیز است.
    \item \textbf{کمترین مربعات متناوب \enparen{Alternating Least Squares - ALS}}: در هر مرحله یکی از ماتریس‌های $P$ یا $Q$ ثابت نگه داشته شده و مسئله به یک دستگاه معادلات خطی مجزا برای هر سطر تبدیل می‌شود. مزیت بی‌نظیر ALS در سیستم‌های توزیع‌شده این است که کاملاً به صورت موازی \enparen{Embarrassingly Parallel} بر روی خوشه‌های MapReduce یا Spark اجرا می‌گردد.
\end{itemize}
\end{examnote}

\begin{summarybox}[جمع‌بندی فصل نهم]
\begin{itemize}
    \item سامانه‌های توصیه‌گر با بهره‌گیری از پدیده دنباله طولانی، گزینه‌های متناسب با سلیقه فردی را پیشنهاد می‌دهند.
    \item فیلترینگ محتوامحور بر شباهت ویژگی‌های کالا تکیه دارد اما دچار محدودیت بیش‌تخصصی شدن است.
    \item فیلترینگ مشارکتی با بهره‌گیری از کسینوس مرکزی، انحرافات مقیاسی کاربران را خنثی می‌نماید.
    \item مدل‌های فاکتورسازی ماتریسی (تجزیه UV) با استخراج ابعاد پنهان و بهینه‌سازی SGD، بالاترین دقت پیش‌بینی را در مقیاس وسیع ارائه می‌دهند.
\end{itemize}
\end{summarybox}
